\section{Related Work}
\label{sec:related-work}

\subsection{Spatial indexes for multidimensional data}

Spatial indexing has a long history, and many of the structures used in modern point-cloud systems come from general multidimensional search. Quadtrees and kd-trees introduced recursive spatial partitioning for composite keys and multidimensional associative search~\cite{finkel1974quadtrees,bentley1975kdtree}. Octrees extend this idea to three-dimensional subdivision and became a natural representation for volumetric and point-based geometry~\cite{meagher1980octree,samet2006foundations}. Region-based structures such as R-trees and R*-trees organize spatial objects by bounding rectangles or boxes and use split heuristics to control overlap, coverage, and occupancy~\cite{guttman1984rtree,beckmann1990rstar}. Bounding-volume hierarchies and linear BVH construction, especially GPU-friendly variants, provide another branch of the same design space, where construction and traversal are shaped by memory layout and parallelism~\cite{karras2012lbvh}.

These structures differ in how they trade off build cost, memory, locality, fanout, overlap, and query selectivity. That observation is central to this work. Instead of treating one primitive as the right answer for all point clouds, our framework treats these primitives as reusable blocks in a measured schema space. The optimizer can still return a single kd-tree, octree, or grid when that is best, but it can also combine primitives when a cloud contains multiple geometric regimes.

\subsection{Point-cloud indexing and processing}

Point-cloud libraries and applications commonly expose kd-tree and octree search because range, radius, and nearest-neighbor queries are basic operations for registration, filtering, reconstruction, segmentation, and visualization~\cite{rusu2011pcl}. Recent point-cloud-specific work continues to refine this space. Some systems improve locality through space-filling curves and linear octrees, reducing cache misses for neighborhood search~\cite{vinambres2026linearoctrees}. Others propose specialized layouts for LiDAR distributions, such as sparse or mixed cell maps designed around the empty-space structure of airborne scans~\cite{laso2025cheesemap}.

Several point-cloud papers already move beyond single static structures. Wang et al. combine an octree with a 3D R*-tree and study the threshold trade-offs between construction cost and nearest-neighbor search~\cite{wang2021hybridpointcloud}. Park et al. use a modifiable nested octree to index BIM surface meshes and point clouds for change detection~\cite{park2021nestedoctree}. Ogayar-Anguita et al. propose nested spatial data structures for LiDAR indexing, allowing hierarchies composed from grids, quadtrees, octrees, and kd-trees~\cite{ogayar2023nestedlidar}. These systems are close in spirit to the present work because they show that hybrid composition is useful for 3D data. They also illustrate the remaining opportunity: the hierarchy is usually designed for a specific application or exposed as a configurable structure, not selected by a repeatable workload-aware search over many candidate schemas.

MultiDataStructure is complementary to these efforts. It does not propose one new fixed point-cloud index. It provides an experimental mechanism for comparing and composing several families under one workload contract. This makes the system closer to an index-design workbench than to a single replacement for existing kd-tree, octree, grid, or BVH implementations.

\subsection{Workload-aware and adaptive physical design}

Database systems have long studied physical design as a problem that depends on the workload. Index selection and physical tuning systems choose auxiliary structures by estimating their effect on a workload, while adaptive indexing methods reorganize data incrementally as queries arrive~\cite{chaudhuri1997indexselection,idreos2007databasecracking,halim2012stochasticcracking}. Recent workload-aware spatial and learned-index systems make the same dependency explicit: WaZI optimizes a learned Z-index for spatial range workloads, WISK learns partitions and hierarchy for spatial keyword queries, and LITune tunes learned-index parameters through reinforcement learning~\cite{pai2023wazi,sheng2023wisk,wang2025litune}. The common idea is that the physical organization should not be separated from the query pattern it is meant to serve.

Our method adopts this workload-aware stance but changes the object being tuned. Rather than selecting secondary indexes for a relational engine, spatial-keyword workload, or one-dimensional learned index, it synthesizes the primary spatial index used by a point-cloud workload. The measured tuple is therefore not a database query plan, but a point cloud, a query profile, and a replayable spatial schema. This distinction is important because the optimizer must account for geometric features such as anisotropy, height distribution, occupancy entropy, and local density, not only logical predicates or table cardinalities.

\subsection{Meta-optimization, synthesis, and learned data structures}

Learned indexes recast some indexing tasks as prediction problems, replacing or augmenting parts of traditional indexes with models that learn the data distribution~\cite{kraska2018learned}. Follow-up systems made learned indexes more tunable, dynamic, and multidimensional. FITing-Tree exposes an error bound that trades space for lookup cost~\cite{galakatos2019fitingtree}; ALEX adapts learned indexes to update-heavy workloads~\cite{ding2020alex}; Flood jointly optimizes multidimensional layout and indexing for range predicates~\cite{nathan2020flood}; and learned spatial indexes apply similar ideas to spatial range queries~\cite{pandey2020learnedspatial}. Recent evolving-index work also treats index shape as something that can be adapted under skewed workloads rather than fixed at creation time~\cite{wang2025shapeshifter}.

A closely related line of work studies data-structure synthesis and meta-optimization. The Data Calculator argues that data structures can be described through fine-grained layout primitives and evaluated through learned cost models before full implementation~\cite{idreos2018datacalculator}. Design-continuum work similarly unifies families such as B+trees, LSM-trees, B-epsilon trees, and hash tables so that a system can reason over a broader design space~\cite{idreos2019learningkv}. In graphics, Herveau et al. show a more domain-specific version of the same idea by combining model prediction and online search to autotune acceleration-structure parameters for ray tracing~\cite{herveau2023autotuning}. These papers are important for the present work because they shift the research question from ``which known data structure should we pick?'' to ``which point in a design space should this workload occupy?''

The present framework follows this meta-optimization viewpoint but differs in two ways. First, it keeps the final decision measurement-driven: candidates are built and queried on the target point cloud, and proxy scores are marked separately from final latency measurements. This is more expensive than pure model-based design-space exploration, but it avoids claiming final speedups from an approximate cost model. Second, it focuses on spatial composition. The learned component is not required to answer queries directly; it may simply rank candidate schemas or prune an expensive search. In this sense, the approach sits between hand-designed spatial data structures and fully learned indexes: it keeps mature spatial primitives, but makes their composition workload-aware and replayable.
